\begin{abstract}
Widrow and Hoff stored neural-network weights in memistors, three-terminal electrochemical cells, in 1960; lithium-ion,
metal-oxide and protonic ECRAM are their modern descendants. We ask, in simulation, how close such devices come to digital
training when every weight update is made of device pulses and the reads are kept ideal (fp32). We compare plain SGD and three
Tiki-Taka-family algorithms on five device specifications, on MNIST, a CIFAR-10 subset and a byte-level language model, with five
seeds that share initialisation and data order across arms. Against SGD with exact updates at the same learning rate,
\RVAllThreeFpN{} pairings (Li-ECRAM and protonic-inspired, each with chopped TTv2 or AGAD) stay within the pre-set margins on all
three tasks, while plain SGD exceeds them on every published device. Against a separately tuned digital Adam no pairing does: on
the language model the best analog result is \RVLmAdamBestGap\,\% worse in bits per byte. Over a factor of four in width no
arm's gap grows detectably beyond the margin; endpoint equivalence was established for one CIFAR pairing only, and growth is detected for another without resolving its size against the margin. A sensitivity sheet with three seeds establishes only a few one-at-a-time
ranges and shows that an earlier asymmetry limit was mostly lost weight range. Against explicit pulse streams the aggregate
update meets preset point-estimate lines on \RYProbeSurMeetsN{} of \RYProbeN{} probes (bootstrap:
\RYProbeSurAtNinetyFiveN), not where clipping removes nearly all of an
update, and ignores shared-line dependence; in training one label and two close rankings change. A conditional energy projection, borrowing an analog inference chip's read efficiency, gives
\RVenSysMin--\RVenChipMax$\times$ less energy per weight-row than an H100, for cells that pass the margin under ideal reads. This
version repairs the simulator after independent reviews and reports how far each repair moved the original results.
\end{abstract}
